\section{Theorem-Level Relation to Neighboring Results}
\label{supp:literature}
We compare the stated versions and risk objects, without claiming an exhaustive exclusion of equivalent consequences. A representation's singular values, a fixed decoder's error, a ridge probe, a cross-context classifier, and the best unregularized population linear risk are different quantities. Nonmonotonicity alone is insufficient to distinguish this work.

\paragraph{Matched functions and joint factor training.}
Kunin et al.~\cite{kunin2024rich}, Sec.~3 and App.~A.2, give exact actual joint gradient flows for supervised $aw^\top x$ models, with conserved imbalance and layer-rate effects. Their Theorems~3.1 and 4.1--4.2 concern implicit bias and effective preconditioning. Our inference is that a rotating feature trajectory could itself yield improvement followed by deterioration on a suitably chosen independent task; absence of a named probe-risk theorem is not an exclusion. The distinction claimed here is the concrete unsupervised square encoder bottleneck, its task-uniform gaps, and its explicit initialization neighborhood.

\paragraph{Same architecture and objective.}
Torralba and Weiss~\cite{torralba2026theory}, Theorem~3.2 and App.~M, analyze convolution, square activation, pooling, and a head under GUPA. Theorem~M.1 assumes equal encoder/head width and a frozen random orthogonal head, and characterizes optimum filter structure. Figures~24--25 already show decreasing objective with deteriorating empirical classification. We do not claim priority for this architecture, loss, factor nonuniqueness, or qualitative degradation.

\paragraph{Layerwise retention, collapse, and stability.}
Xue et al.~\cite{xue2024benefits}, Theorems~3.3--3.6, establish weighting/balance and downstream sample-complexity results; Theorem~3.8 treats actual nonlinear BReLU flow with asymptotic coordinate suppression after projection and retention before it. Their App.~D.2 also reports early-stopping benefits. Tepe~\cite{tepe2024dynamics}, Theorem~2 and Propositions~1--2, studies spectral-contrastive Gram convergence and deep-linear singular-value bounds under commutation or directional assumptions. Liao et al.~\cite{liao2026imbalanced}, Theorems~3.1--3.2, analyzes finite stages and pruning for imbalanced contrastive training with stop-gradient and weight decay; its population risk is pretraining loss. Jing et al.~\cite{jing2022collapse} proves alignment/collapse results for linear contrastive dynamics. Tian~\cite{tian2023nonlinearity}, Lemma~4 and Theorems~3,5--6, includes nonlinear mode dynamics and stability; homogeneity permits monomials, so it cannot be dismissed as excluding square nonlinearities. None of these comparisons makes singular-value decay alone evidence of a finite-time, preserved-rank population OLS risk gap.

\paragraph{Finite geometric overshoot is an existing result.}
Yang, Saxe, and Latham~\cite{yang2026abstraction}, Assumption~1, optimize out a ridge readout at every time. Proposition~B.2 proves finite-time abstraction overshoot; Theorem~5 analyzes initialization-scale dependence of its peak; Proposition~6 gives near-perfect-alignment residence of order $\log(1/\kappa)$. Proposition~1 relates abstraction to a cross-context classification probe under additional distributional conditions. Their exact reduced dynamics therefore supplies a close nonmonotonic-geometry and duration precedent. Our risk reoptimizes an unregularized population decoder on the whole distribution at each state, and the pretraining head follows actual finite-rate joint dynamics. These distinctions, not generic overshoot, define the comparison.

\paragraph{Quadratic moment closure is also established.}
Neural quadratic forms~\cite{liu2026nqf}, Theorem~2, closes actual SGD on first and second parameter moments, retaining the finite-step quadratic term. Its Theorem~7 gives an isotropic matrix-Riccati solution and commuting modal reductions for supervised square loss. Yang~\cite{yang2026spectral}, Eq.~(1) and Theorems~4.1,4.3, studies factorized centered quadratic learners with a fixed positive teacher, proving between-distribution loss contrasts and their coupling order. These establish close tools based on matrix products and moment geometry. The trainable-head square model is cubic in its joint parameters; absorbing positive heads into quadratic filters changes their Euclidean flow. Neither citation is used as an exhaustive negative claim about what could be deduced from its results.

The central result locates readability loss in the encoder itself during actual unsupervised joint training. Uniform finite gaps over a prespecified task cone show a shared finite-amplitude effect. Preserved rank identifies reallocation among readable directions, and optimal population probes separate it from finite-label estimation and invertible feature scaling. The explicit six-coordinate neighborhood establishes persistence beyond the symmetric solution. These properties specify the scientific consequence of the construction and its local stability guarantee.
